\section*{Chapter \ref{ch:ll:cpu-microarchitecture} --- CPU Microarchitecture}

\begin{solution}{exo:ll:cpu-microarchitecture:1}
$600 \times 4.0 = 2\,400$ cycles, so $2\,400 \times 2.4 = 5\,760$ instructions.
\end{solution}

\begin{solution}{exo:ll:cpu-microarchitecture:2}
$3\,000 / 3.0 = \qty{1\,000}{\nano\second}$, which is $4\,500$ core cycles at 4.5~GHz. The TSC counts time.
\end{solution}

\begin{solution}{exo:ll:cpu-microarchitecture:3}
$4 \times 2 = 8$ chains (\cref{prop:ll:cpu-microarchitecture:little}); three chains complete $3/4$ of an operation a
cycle, $37.5\%$ of the throughput.
\end{solution}

\begin{solution}{exo:ll:cpu-microarchitecture:4}
Branchy: $0.3 \times 20 = 6$ cycles on average; branchless: 4. The branchless version is faster by 2 cycles a message
(\cref{prop:ll:cpu-microarchitecture:breakeven}: $m = 0.3 > 4/20$).
\end{solution}

\begin{solution}{exo:ll:cpu-microarchitecture:5}
On the committed measurement, $(2.80 - 0.21)/0.5 = \qty{5.2}{\nano\second}$, about 24 cycles at the measured
4.6~GHz of CPU~2. The estimate assumes that the predictor mispredicts half the iterations on shuffled data (it can do
no better than chance on independent random values) and that nothing else differs between the two runs.
\end{solution}

\begin{solution}{exo:ll:cpu-microarchitecture:6}
Each instruction of a chain needs the previous one's result, so at most one is in flight however many slots the buffer
has; the chain runs at its latency. The buffer hides latency by overlapping the chain with independent work (other
chains, the next iterations' loads), and hides cache misses as long as that work fits in it.
\end{solution}

\begin{solution}{exo:ll:cpu-microarchitecture:7}
\texttt{measured\_blocks.csv}: blocks of 1, 2, 4 and 8 run at the sorted speed (the predictor, which uses the recent
history of the branch, learns patterns of that period completely); blocks of 16 and 64 cost about twice as much, a
misprediction at each switch; blocks of 1\,024 are back at the sorted speed, since switches are rare.
\end{solution}

\begin{solution}{exo:ll:cpu-microarchitecture:8}
Three flaws. \texttt{add \$1} chains are resolved in the renamer, about five a cycle on this core family, so the
``cycle'' reference is five times too short. A warm core running one order ten million times keeps its caches and
predictor perfectly trained, which the hot path never sees after a quiet period. And the counter must be fenced and
its overhead subtracted. Use a register add chain for the clock, vary the orders, and measure cold as well as warm.
\end{solution}

\begin{solution}{pb:ll:cpu-microarchitecture:1}
Figures from the committed measurement (\cref{fig:ll:cpu-microarchitecture:branch,fig:ll:cpu-microarchitecture:freq}).
\begin{enumerate}
\item Sorted: branchy \qty{0.21}{\nano\second}, branchless \qty{0.32}{\nano\second} per element; shuffled: branchy
\qty{2.80}{\nano\second}, branchless \qty{0.32}{\nano\second}.
\item The predictor is almost always right, and the branchy loop does no work for the half of the elements below 128;
the branchless loop computes and adds a mask for every element.
\item $(2.80 - 0.21)/0.5 = \qty{5.2}{\nano\second}$.
\item $5.2 \times 4.64 \approx 24$ cycles.
\item Branchy: $0.21 + 5.2\,m$ ns; branchless: \qty{0.32}{\nano\second}.
\item $m^\ast = 0.11/5.2 \approx 2.1\%$.
\item Branchy, $0.21 + 0.05 = \qty{0.26}{\nano\second}$ against 0.32: by about \qty{0.06}{\nano\second} a trade.
\item Branchless: $0.21 + 0.35 \times 5.2 = \qty{2.03}{\nano\second}$ against about 0.32: by about
\qty{1.7}{\nano\second} a trade.
\item $50\,000 \times 0.35 \times \qty{5.2}{\nano\second} \approx \qty{91}{\micro\second}$ a second.
\item Small: \qty{5}{\nano\second} is about 0.4\% of \qty{1.3}{\micro\second}; the cost is in the clustering,
not the average.
\item Busy markets are the ones where the side of the next trade is least predictable, and they are when latency is
worth most.
\item It would make the branchy version cheap again; but the pattern of market sides is not periodic, so there is little
to learn.
\item \textbf{Named result.} A misprediction costs about \qty{5}{\nano\second} on this laptop, some 24 cycles; the
branchless rewrite wins once more than about 2\% of the branches mispredict.
\item Virtualisation, a clock below turbo under load, and a benchmark that also pays for the loop's other work; the
published figure is a minimum on bare hardware.
\item Elapsed time at a constant rate (about 3~GHz here), not the core's cycles, which follow its clock.
\item To keep the thread on one core with its caches and predictor; under a hypervisor it pins to a virtual CPU,
which may move between physical cores.
\item Likely turn the branch into a conditional move or vectorise the loop, removing the misprediction and the
experiment.
\item Predictable: the message-type dispatch on a feed dominated by one type. Unpredictable: whether a trade was at the
bid or the ask.
\item With hardware performance counters (\texttt{perf stat -e branch-misses}) on the production host, or by counting
the side changes in recorded market data.
\item When a branch depends on unpredictable data and its misprediction rate exceeds the ratio of the rewrite's extra
cost to the penalty.
\end{enumerate}
\end{solution}

\begin{solution}{iq:ll:cpu-microarchitecture:1}
\omterm{def:ll:cpu-microarchitecture:branch}{Branch prediction}. On sorted data the \texttt{if} goes the same way for long runs and is predicted correctly; on random
data it is wrong half the time, and each misprediction flushes the pipeline, costing tens of cycles.
\iqlookfor{the pipeline flush on a wrong guess, and that the work itself did not change.}
\end{solution}

\begin{solution}{iq:ll:cpu-microarchitecture:2}
The core starts each instruction when its inputs are ready, not in program order, and retires results in order through
the \omterm{def:ll:cpu-microarchitecture:ooo}{reorder buffer}. It is limited by dependency chains, by the size of the buffer and scheduler, by \omterm{def:ll:cpu-microarchitecture:branch}{branch
mispredictions} (which discard the work) and by the width of the front end.
\iqlookfor{readiness-driven scheduling with in-order retirement, and the limits.}
\end{solution}

\begin{solution}{iq:ll:cpu-microarchitecture:3}
One accumulator is a dependency chain running at the latency of an addition; four chains fill the adders, whose
throughput is latency times width (Little's law). The compiler may not reassociate floating-point additions, which
changes the result, unless told it may (\texttt{-ffast-math}, or \texttt{-fassociative-math}).
\iqlookfor{latency versus throughput and floating-point non-associativity.}
\end{solution}

\begin{solution}{iq:ll:cpu-microarchitecture:4}
Call it many times between fenced reads of the \omterm{def:ll:cpu-microarchitecture:tsc}{time-stamp counter} (or time a batch), keep its result alive with a
\omterm{def:ll:cpu-microarchitecture:ubench}{compiler barrier}, subtract the timer's overhead, pin the thread, warm up, and report percentiles. Pitfalls: the
optimiser deleting or hoisting the work, a trained predictor and warm caches, frequency changes, TSC ticks mistaken for
cycles.
\iqlookfor{barriers, fences, overhead, and the difference between warm and realistic conditions.}
\end{solution}

\begin{solution}{iq:ll:cpu-microarchitecture:5}
No: the sibling shares the execution units, first-level caches and predictors, so its load adds latency and \omterm{def:ll:where-latency-comes-from:tail}{jitter}.
Leave it idle or disable \omterm{def:ll:cpu-microarchitecture:smt}{simultaneous multithreading} on the hot cores.
\iqlookfor{shared per-core resources and the \omterm{def:ll:where-latency-comes-from:tail}{jitter} they create.}
\end{solution}

\begin{solution}{iq:ll:cpu-microarchitecture:6}
The core woke from a deep \omterm{def:ll:cpu-microarchitecture:dvfs}{idle state} (limit C-states and compare); it runs at a low clock before scaling up (fix the
frequency and compare); caches, translation buffers and predictors were cold (keep the path warm with synthetic
traffic and compare). Confirm each by timing the first message with and without the fix.
\iqlookfor{\omterm{def:ll:cpu-microarchitecture:dvfs}{idle states}, frequency scaling, and cold microarchitectural state, each with an experiment.}
\end{solution}
